\subsectionaddtolist{۴.۱ تکامل مدل اعطای مجوز}

\subsubsection*{آ}
\begin{itemize}
    \item {فهرست نصب (\lr{Install-time Manifest}):} 
    \begin{itemize}
        \item \textit{ضعف امنیتی:} مشکل همه یا هیچ و پذیرش کورکورانه. کاربر برای نصب برنامه مجبور است تمام مجوزها را یکجا تایید کند. از آنجا که کاربران معمولا متن طولانی مجوزها را نمی‌خوانند، برنامه‌های مخرب به راحتی به همه چیز دسترسی دائمی و پس‌زمینه پیدا می‌کنند.
        \item \textit{ضعف کاربردپذیری:} کاربر در لحظه نصب هیچ کانتکستی ندارد که بداند چرا برنامه مثلا به دوربین نیاز دارد، چون هنوز با محیط برنامه و فیچرهای آن کار نکرده است.
    \end{itemize}
    \item {درخواست لحظه‌ای (\lr{Time-of-use Prompt}):}
    \begin{itemize}
        \item \textit{ضعف امنیتی:} خستگی از هشدارها. اگر برنامه پشت سر هم پاپ‌آپ نشان دهد، کاربر برای رد شدن از آن بدون دقت تایید را می‌زند. همچنین در مدل‌های اولیه، وقتی کاربر در لحظه تایید را می‌زد، مجوز برای همیشه (حتی در پس‌زمینه) به برنامه داده می‌شد.
        \item \textit{ضعف کاربردپذیری:} قطع کردن جریان کار کاربر. اگر پیام درست در لحظه نامناسبی ظاهر شود، باعث آزار کاربر شده و ممکن است بدون فکر آن را رد کند که باعث از کار افتادن ویژگی‌های برنامه می‌شود.
    \end{itemize}
\end{itemize}

\subsubsection*{ب}
\begin{itemize}
    \item {ضعفی که برطرف شد:} مشکل فقدان کانتکست و پذیرش کورکورانه در زمان نصب حل شد. کاربر حالا فقط زمانی پیام را می‌بیند که مثلا روی دکمه عکس‌گرفتن کلیک کرده است، پس منطق درخواست مجوز را درک می‌کند و دسترسی‌ها به صورت \lr{Granular} اعطا می‌شوند.
    \item {ضعفی که باقی ماند:} مشکل دسترسی دائمی پس از یک بار تایید. اندروید ۶ هنوز تفاوت بین دسترسی موقت و دسترسی دائمی را به خوبی تفکیک نکرده بود؛ برنامه‌ای که یک بار برای یک کار خاص مجوز می‌گرفت، می‌توانست در پس‌زمینه هم به دسترسی ادامه دهد. همچنین خطر خستگی از پاپ‌آپ‌ها به قوت خود باقی ماند.
\end{itemize}

\subsubsection*{ج}
\begin{itemize}
    \item {مجوز یک‌بار مصرف (\lr{Only this time}):} دقیقا کاستی دسترسی دائمی پس از تایید را هدف گرفته است. با این کار، برنامه فقط در همان نشست و تا زمانی که باز است مجوز دارد و برای دفعه بعد باید دوباره اجازه بگیرد. این یعنی جلوگیری از سوءاستفاده‌های پس‌زمینه.
    \item {موقعیت تقریبی در برابر دقیق:} کاستی بیش‌امتیازی را هدف می‌گیرد. بسیاری از برنامه‌ها (مثل هواشناسی) فقط نیاز دارند بدانند کاربر در کدام شهر است، نه مختصات دقیق \lr{GPS}. این ویژگی به کاربر حق انتخاب برای اعمال اصل حداقل دسترسی را می‌دهد.
    \item {شفافیت ردیابی برنامه:} کاستی ردیابی نامرئی و تبانی شرکت‌های تبلیغاتی را هدف قرار داده است. برنامه‌ها با استفاده از شناسه‌های یکتا (مثل \lr{IDFA}) داده‌های رفتاری کاربر را به اشتراک می‌گذاشتند. این سازوکار، ردیابی بین‌برنامه‌ای را از حالت پیش‌فرض پنهان، به حالت شفاف نیازمند تایید (\lr{Opt-in}) تبدیل کرد.
\end{itemize}

\subsectionaddtolist{۴.۲ دور زدن کامل سامانه مجوز}

\subsubsection*{آ}
\begin{itemize}
    \item {کانال جانبی:} زمانی اتفاق می‌افتد که یک برنامه به تنهایی و بدون داشتن مجوز یک داده حساس، از طریق یک مسیر یا منبع محافظت‌نشده (مثلا فراداده‌های یک فایل دیگر) به آن داده دسترسی پیدا می‌کند. سیستم‌عامل در محافظت از مرزهای داده ناکام بوده است.
    \item {کانال پنهان:} زمانی رخ می‌دهد که دو برنامه مجزا روی دستگاه با هم تبانی می‌کنند. برنامه‌ای که دسترسی دارد داده را می‌خواند و از طریق یک منبع مشترک سیستمی ظاهرا بی‌خطر، آن را به برنامه‌ای که مجوز شبکه دارد می‌رساند.
\end{itemize}
هر دو روش اصل ایزوله‌سازی برنامه‌ها را دور می‌زنند؛ اما کانال جانبی از ضعف خود سیستم در نشت اطلاعات بهره می‌برد، در حالی که کانال پنهان از تبانی اپلیکیشن‌ها و دور زدن مرزهای ارتباطی استفاده می‌کند.

\subsubsection*{ب}
\begin{itemize}
    \item {نمونه اول (استخراج موقعیت از روی \lr{EXIF} عکس):} یک {کانال جانبی} است. داده هدف موقعیت مکانی بوده که از مسیر فراداده‌های تصاویر گالری به دست آمده است.
    \item {نمونه دوم (تخمین موقعیت از روی \lr{BSSID} و وای‌فای):} یک {کانال جانبی} است. داده هدف موقعیت مکانی بوده که از مسیر اطلاعات شبکه وای‌فای و شناسه روترها و نگاشت آن‌ها استخراج شده است.
    \item {نمونه سوم (نوشتن \lr{IMEI} روی حافظه مشترک):} یک {کانال پنهان} (تبانی) است. داده هدف شناسه دائمی دستگاه بوده که توسط برنامه مجاز روی حافظه مشترک نوشته شده تا برنامه بدون مجوز آن را بخواند.
\end{itemize}

\subsubsection*{ج}
اصلاح پنجره درخواست مجوز فایده‌ای ندارد، چون برنامه اصلا از آن مسیر عبور نمی‌کند. کانال‌های جانبی از {سطح حمله غیرمنتظره} استفاده می‌کنند؛ طراح سیستم‌عامل در نظر نگرفته که مثلا از روی فراداده عکس، وضعیت شبکه وای‌فای، یا دمای باتری بشود موقعیت یا اطلاعات حساس کاربر را بیرون کشید. چون این مسیرها محافظت‌شده نبودند، شناسایی و بستن تمام این روزنه‌ها بدون خراب کردن کارکرد عادی سیستم‌عامل بسیار دشوار است.

\subsubsection*{د}
\begin{itemize}
    \item {پروتکل و سازوکار کدگذاری:} منبع مشترک انتخابی را حجم اشغال‌شده حافظه نهان در نظر می‌گیریم.
    \begin{itemize}
        \item {کدگذاری:} برنامه A برای فرستادن بیت \texttt{1} حجم مشخصی از حافظه نهان را پر می‌کند و برای فرستادن بیت \texttt{0} آن را خالی می‌کند. برنامه B با چک کردن حجم اشغال‌شده، صفر یا یک را می‌خواند.
        \item {همگام‌سازی:} زمان را به بازه‌های ۵۰ میلی‌ثانیه‌ای تقسیم می‌کنیم. در ۲۰ میلی‌ثانیه اول، A تغییر را اعمال می‌کند و در ۲۰ میلی‌ثانیه دوم، B وضعیت را می‌خواند.
    \end{itemize}
    \item {نرخ نشت داده و زمان انتقال:} با فرض بازه ۵۰ میلی‌ثانیه‌ای، نرخ انتقال {۲۰ بیت بر ثانیه} خواهد بود. برای انتقال یک شماره تلفن ۵۰ بیتی، تنها {{5.2} ثانیه} زمان نیاز است.
    \item {عوامل محدودکننده و اهمیت تهدید:} در عمل، نویز سیستم (تغییر حجم حافظه توسط سایر برنامه‌ها) و تاخیر زمان‌بندی پردازنده این نرخ را محدود می‌کنند. با این حال، حتی یک نرخ بسیار پایین (مثلا چند بیت در ثانیه) یک تهدید بسیار جدی است؛ زیرا داده‌های حیاتی مثل رمز عبور، شماره کارت یا IMEI حجم بسیار کمی دارند و به تدریج در خفا به طور کامل لو می‌روند.
\end{itemize}